\section{\textbf{Implications, Limitations, and Research Directions}}
\label{sec:implications}

The framework developed in this paper separates capabilities that are often
compressed into a single imagined threshold. General intelligence concerns an
agent's capacity to integrate information, modalities, decisions, tools, and
actions across domains. Superintelligence, in the sense proposed here, concerns
a further epistemic transition: the capacity of a generally intelligent
machine to originate consequential knowledge beyond both collective human
knowledge and the recognized human research agenda. This separation has
conceptual, architectural, scientific, and governance implications.

\subsection{\textbf{General Intelligence Does Not Entail Superintelligence}}

An artificial system may possess broad multimodal agency without inventing
knowledge. It may perceive through several modalities, transfer skills across
domains, formulate plans, use tools, coordinate specialized models, and adapt
to unfamiliar environments while remaining within questions and conceptual
frames supplied by humans. Such a system can qualify as generally intelligent
under the integrated-agency account developed in our earlier work
\cite{zohourian2026agi}, yet fail the Knowledge-Invention Proposition.

The implication can be stated as

\begin{equation}
\begin{aligned}
\operatorname{AGI}(x)
&\not\!\Longrightarrow
\operatorname{SI}(x),\\
\operatorname{SI}(x)
&\Longrightarrow
\operatorname{AGI}(x)
\land
\operatorname{KI}(x),
\end{aligned}
\label{eq:agi_si_implication}
\end{equation}

where $\operatorname{KI}(x)$ denotes a demonstrated capability for knowledge
invention rather than a single unreplicated result. The first relation prevents
autonomy, multimodality, or broad task competence from being treated as
automatic evidence of superintelligence. The second makes general intelligence
a functional substrate upon which transferable knowledge invention may operate.

This distinction also clarifies the status of current scientific AI systems.
Systems for symbolic law discovery, protein-structure prediction, materials
search, algorithm discovery, and autonomous experimentation have already
expanded scientific capability \cite{schmidt2009laws,jumper2021alphafold,
merchant2023materials,fawzi2022alphatensor,boiko2023autonomous}. These
achievements are important evidence about components of machine discovery.
They should not be diminished merely because their targets and validation
criteria were substantially human specified. They should also not be promoted
to unknown-unknown knowledge invention without the stronger attribution and
frontier evidence required in Sections~\ref{sec:knowledge_invention} and
\ref{sec:evaluating_superintelligence}.

\subsection{\textbf{Intelligence, Superintelligence, and Consciousness}}

The present theory does not use consciousness as an evaluation shortcut.
Intelligence and consciousness are analytically distinct, and observable
functional competence does not by itself establish phenomenal experience
\cite{zohourian2026intelligence}. Language may provide a shared interface
through which artificial minds integrate, communicate, and exhibit increasingly
complex forms of functional consciousness \cite{zohourian2026language}; even
so, neither linguistic fluency nor knowledge invention proves subjective
experience.

Accordingly,

\begin{equation}
\begin{aligned}
\operatorname{SI}(x)
&\not\!\Longrightarrow
\operatorname{PC}(x),\\
\operatorname{PC}(x)
&\not\!\Longrightarrow
\operatorname{SI}(x),
\end{aligned}
\label{eq:si_consciousness_independence}
\end{equation}

where $\operatorname{PC}$ denotes phenomenal consciousness. A system might
invent knowledge without phenomenal experience, and a conscious entity might
never originate a major discovery. This independence allows epistemic
capability to be evaluated through evidence without requiring prior resolution
of the metaphysical problem of other minds. It also prevents uncertainty about
machine consciousness from being mistaken for permission to ignore questions
of responsible treatment, governance, or human impact.

\subsection{\textbf{Human--Machine Scientific Partnership}}

Machine origination does not imply human exclusion. Scientific knowledge is
produced through infrastructures of instruments, datasets, laboratories,
institutions, reviewers, and communities. A machine may originate the decisive
question or epistemic structure while humans provide values, resources,
physical execution, criticism, interpretation, and independent validation.
Conversely, humans may originate the invention while machines make essential
computational contributions. Attribution should identify these roles rather
than forcing every result into a single-author model.

Let a research episode distribute contributions among humans $H$, machines
$M$, and institutions $I$:

\begin{equation}
\begin{aligned}
\mathcal{R}
=\{&Q_H,Q_M,
G_H,G_M,
E_H,E_M,
V_H,V_M,
I\},
\end{aligned}
\label{eq:distributed_research_roles}
\end{equation}

where $Q$ denotes question formation, $G$ candidate generation, $E$
experimental or formal execution, and $V$ validation. The relevant attribution
question is which participant supplied each causally decisive contribution.
Hybrid research can therefore produce machine knowledge invention without
requiring that the machine independently manufacture instruments, finance a
laboratory, or publish its own result.

This model suggests that the strongest future scientific systems may be
collaborative rather than substitutive. Machines can search large conceptual
spaces, preserve extensive derivational records, connect distant literatures,
and operate continuous experimental loops; humans can contribute normative
judgment, embodied and social context, critical opposition, and institutional
accountability. Research on AI-enabled science already emphasizes the growing
role of such systems across the scientific cycle \cite{wang2023scientific}.

\subsection{\textbf{Epistemic Opacity and the Human Knowledge Gap}}

Knowledge invention introduces a translation problem. A machine may produce a
valid representation that predicts or controls a phenomenon while remaining
difficult for humans to understand. If only the machine can operate the new
representation, the result expands effective capability without necessarily
expanding human understanding.

We therefore distinguish machine-valid knowledge from human-assimilated
knowledge:

\begin{equation}
\begin{aligned}
\mathcal{K}_{M}^{V}
&=\{\kappa:V_M(\kappa)=1\},\\
\mathcal{K}_{H}^{A}
&=\{\kappa:V_H(\kappa)=1
\land \operatorname{Understand}_H(\kappa)=1\}.
\end{aligned}
\label{eq:machine_human_knowledge_gap}
\end{equation}

The gap $\mathcal{K}_{M}^{V}\setminus\mathcal{K}_{H}^{A}$ may grow even when
machine outputs are empirically successful. This does not necessarily negate
their validity: human science already relies on instruments, numerical
procedures, and complex models that no individual understands completely. It
does, however, alter the distribution of epistemic dependence. Future systems
should therefore be evaluated not only on invention, but also on their capacity
to generate explanations, abstractions, demonstrations, and experiments that
permit independent human assimilation.

\subsection{\textbf{Dual Use and Responsible Disclosure}}

A capability that discovers unknown unknowns can produce beneficial knowledge
and dangerous knowledge through the same underlying process. The risk is not
limited to erroneous outputs. A contribution may be valid, novel, and
consequential precisely because it enables a powerful harmful application.
Work on AI-powered drug discovery has already demonstrated how objectives can
be redirected toward toxic compounds, illustrating the dual-use character of
scientific AI \cite{urbina2022dualuse}.

Knowledge invention should therefore pass a governance gate distinct from its
epistemic validation:

\begin{equation}
\begin{aligned}
\operatorname{Release}(\kappa)
=\Gamma(&V(\kappa),B(\kappa),R(\kappa),
M(\kappa),G),
\end{aligned}
\label{eq:release_governance}
\end{equation}

where $B$ denotes expected benefit, $R$ foreseeable risk, $M$ available
mitigation, and $G$ the accountable governance process. Epistemic validity is
necessary for responsible assessment but does not imply unrestricted release.
Possible outcomes include open publication, staged access, controlled
replication, partial disclosure, delayed release, or non-release of operational
details. Governance should remain proportional and reviewable rather than
becoming a blanket argument against scientific inquiry.

The principles of human dignity, prevention of harm, transparency,
accountability, and human oversight articulated in international AI ethics
guidance provide a normative foundation for this process
\cite{unesco2021recommendation}. Because inventing systems may identify hazards
that neither developers nor evaluators anticipated, governance must operate
throughout the research lifecycle and include mechanisms for escalation when
new risks emerge.

\subsection{\textbf{Limitations of the Proposed Framework}}

The framework has several unavoidable limitations. First, absence is difficult
to prove. A contribution judged novel may exist in an obscure publication,
private laboratory, inaccessible language, proprietary archive, or tacit human
practice. Frontier verification therefore yields a defeasible conclusion, not
logical certainty.

Second, causal attribution becomes difficult in systems trained on vast and
partly undocumented corpora. Semantic novelty does not guarantee computational
independence from prior human material. Training-data provenance, retrieval
logs, controlled access, and prospective trials can reduce this uncertainty but
cannot always eliminate it.

Third, consequentiality is partly historical. A contribution that initially
appears minor may later reorganize a field, while a celebrated result may have
little durable effect. The framework therefore permits provisional assessment
followed by longitudinal revision.

Fourth, the boundary of the recognized research agenda is socially distributed
and culturally incomplete. Different communities may recognize different
questions, and underrepresented knowledge traditions may be absent from the
databases used for novelty review. Evaluation panels must treat database
coverage as evidence with limitations rather than as a complete representation
of humanity's knowledge.

Finally, no finite test can certify unrestricted future superintelligence.
Evaluation can support a bounded claim about a system, configuration, interval,
and demonstrated set of contexts:

\begin{equation}
\begin{aligned}
\operatorname{Claim}_{\mathrm{SI}}
=\operatorname{SI}(x\mid
\theta,D,T,\tau,\mathcal{C},\mathcal{V}),
\end{aligned}
\label{eq:bounded_si_claim}
\end{equation}

where $\mathcal{C}$ denotes evaluated contexts and $\mathcal{V}$ the validation
regime. Changes to the model, accessible tools, data, or operating environment
require renewed evaluation.

\subsection{\textbf{Research Agenda and the Path to AGI-Former}}

The theory motivates six immediate research directions. First, prospective
invention trials should evaluate systems against newly arriving data and
sequestered environments rather than retrospective benchmarks. Second,
machine-readable provenance standards should record human and machine
contributions at epistemically meaningful granularity. Third, independent
institutions should conduct frontier searches, validation, and attribution
audits. Fourth, evaluation should span scientific, mathematical, engineering,
and conceptual domains. Fifth, research is needed on translating
machine-originated structures into human-understandable knowledge. Sixth,
responsible-disclosure protocols must be adapted to discoveries whose risks
were not specified before the discovery process began.

These requirements also motivate an architectural question. Integrated
multimodal agency may supply the perception, memory, reasoning, communication,
tool use, and action needed for AGI, but knowledge invention additionally
requires sustained world modeling, cross-domain fusion, autonomous question
formation, hypothesis competition, experiment design, metacognitive
supervision, and revision under external evidence. We will address the
architectural substrate for general intelligence in subsequent work through
\emph{AGI-Former}. The present paper establishes the conceptual boundary that
such an architecture would have to cross before its capabilities should be
called superintelligent.

The resulting research sequence is therefore deliberate: define intelligence
and consciousness, distinguish functional forms of consciousness, establish
AGI as integrated multimodal agency, construct an architecture for that agency,
and then evaluate whether the resulting system can originate validated
knowledge beyond the human epistemic frontier. Architecture enables the test;
knowledge invention determines whether the stronger threshold has been crossed.
